\chapter{Alpha-1 Antitrypsin Deficiency}

\section*{Definition and Overview}
Alpha-1 antitrypsin deficiency (AATD) is an inherited condition in which the body produces insufficient levels of alpha-1 antitrypsin (AAT), a protective protein produced chiefly by the liver. In health, AAT neutralises proteases released by white blood cells during inflammation, shielding the lungs from tissue damage. In AATD, reduced protection leaves the lung vulnerable to progressive destruction, producing early-onset emphysema, particularly in smokers. In addition, the abnormal AAT protein folds incorrectly and accumulates within liver cells, causing liver injury that ranges from abnormal liver tests to cirrhosis. AATD is one of the more common inherited conditions in people of European descent, yet it remains substantially under-diagnosed, with many affected individuals identified only after years of respiratory symptoms.

\section*{Epidemiology}
AATD is estimated to affect roughly 1 in 2,000 to 1 in 5,000 people in many countries, although the prevalence varies widely by population and ethnic background. Many more people carry a single copy of the faulty gene without developing disease, but they may still be at increased risk of lung problems, especially if they smoke. Men and women are affected equally. Lung disease typically becomes apparent between the ages of 30 and 50 in people who smoke, and later, or not at all, in non-smokers. A significant minority of affected individuals develop liver problems, and liver disease can present in childhood as neonatal hepatitis, with most affected children recovering, although a proportion develop cirrhosis over time.

\section*{Aetiology and Pathophysiology}
AATD is caused by mutations in the SERPINA1 gene, which provides instructions for making alpha-1 antitrypsin. The condition is inherited in an autosomal co-dominant pattern, meaning that the severity depends on the combination of gene variants (alleles) inherited from both parents. The most common clinically important variant is the Z allele; a person who inherits two Z alleles (PiZZ) has severe deficiency, while those with one Z and one normal allele (PiMZ) have intermediate levels and generally milder disease.

The two main disease mechanisms are:

\begin{itemize}
    \item \textbf{Lung disease:} low circulating AAT leaves lung tissue inadequately protected. When neutrophils release proteases such as neutrophil elastase during infection or inflammation, the enzyme is free to destroy the elastic tissue of the air sacs, producing emphysema. Smoking both increases the number of neutrophils in the lung and further inactivates AAT, greatly accelerating damage.
    \item \textbf{Liver disease:} the mutant Z protein folds incorrectly, forms polymers, and accumulates inside hepatocytes, where it causes injury and progressive scarring (cirrhosis) independently of smoking or alcohol.
    \item \textbf{Other associations:} panniculitis, a painful inflammatory skin condition, and vasculitis affecting small vessels are recognised extra-pulmonary complications of severe deficiency.
\end{itemize}

\section*{Clinical Features}
Symptoms often resemble ordinary lung disease, which is why the condition is frequently missed:

\begin{itemize}
    \item \textbf{Lungs:} progressive shortness of breath, particularly on exertion; persistent cough with or without phlegm; wheeze; and frequent or prolonged chest infections.
    \item \textbf{Distribution of emphysema:} involvement tends to be more prominent at the bases of the lungs, in contrast to the upper-lobe predominance of typical smoking-related emphysema.
    \item \textbf{Liver:} in some people, yellowing of the skin or eyes (jaundice), swelling of the abdomen or legs, and abnormal liver blood tests; infants may present with prolonged neonatal jaundice.
    \item \textbf{Other:} fatigue, reduced exercise tolerance, and a family history of early emphysema or unexplained liver disease in close relatives.
\end{itemize}

A striking feature of AATD is the onset of respiratory symptoms at a relatively young age, often before 45, particularly in those who smoke.

\section*{Differential Diagnosis}
Conditions that can present similarly include:

\begin{itemize}
    \item Chronic obstructive pulmonary disease due to smoking without AATD.
    \item Asthma, particularly when wheeze and breathlessness are prominent.
    \item Chronic bronchitis and recurrent respiratory infections.
    \item Emphysema from other causes, such as severe tobacco use or occupational dust exposure.
    \item Other inherited lung conditions, including cystic fibrosis.
    \item Fatty liver disease, viral hepatitis, and other causes of abnormal liver tests.
\end{itemize}

The diagnosis of AATD should be considered in any adult with emphysema, especially if onset is young, the pattern is atypical, or there is a family history.

\section*{When to Seek Medical Care}
See a doctor if you have persistent breathlessness, cough, wheeze, or repeated chest infections, particularly if symptoms began at a relatively young age, if you have a family history of early emphysema or unexplained liver disease, or if liver tests have been abnormal without an obvious cause.

\textbf{Call 000 (or local emergency services) for:}
\begin{itemize}
    \item Sudden severe breathlessness or chest pain.
    \item Coughing up blood.
    \item Blue lips or confusion, which may indicate dangerously low oxygen levels.
    \item Collapse or loss of consciousness.
\end{itemize}

\section*{Diagnosis and Investigations}
AATD is confirmed by simple blood tests, so a high index of suspicion is important:

\begin{itemize}
    \item \textbf{Blood test:} measurement of the serum alpha-1 antitrypsin level, which is low in significant deficiency.
    \item \textbf{Genetic testing (genotyping):} identification of the exact gene variants, usually from a blood or cheek swab sample, which defines the pattern of inheritance and risk.
    \item \textbf{Lung function tests (spirometry):} measurement of airflow limitation and lung volumes to detect emphysema.
    \item \textbf{Chest imaging:} chest X-ray and high-resolution CT to demonstrate the pattern and severity of emphysema.
    \item \textbf{Liver tests:} blood tests of liver enzymes and function, with ultrasound and sometimes liver biopsy or elastography to assess fibrosis.
    \item \textbf{Family screening:} offering testing to close relatives, because early detection allows preventive measures that improve outlook.
\end{itemize}

\section*{Management}
There is no cure for AATD, but treatment slows progression and controls symptoms:

\begin{itemize}
    \item \textbf{Smoking cessation is the single most important step}, together with avoiding second-hand smoke, occupational dusts, and fumes.
    \item \textbf{Inhaler therapy:} bronchodilators and inhaled corticosteroids, as used in COPD and asthma, to relieve breathlessness and reduce exacerbations.
    \item \textbf{Vaccination:} influenza and pneumococcal vaccines to reduce the risk of serious chest infections.
    \item \textbf{Pulmonary rehabilitation:} a supervised program of exercise training and education that improves symptoms and quality of life.
    \item \textbf{Augmentation therapy:} regular intravenous infusions of pooled human AAT, offered in specialist centres to selected non-smoking individuals with established emphysema.
    \item \textbf{Oxygen therapy:} for people with persistent low blood oxygen levels.
    \item \textbf{Liver care:} monitoring of liver tests, a healthy lifestyle, avoidance of excessive alcohol, and referral for transplant assessment in advanced cirrhosis.
    \item \textbf{Regular follow-up:} with a respiratory physician and, where indicated, a hepatologist.
\end{itemize}

\section*{Complications}
AATD can lead to serious complications:

\begin{itemize}
    \item Progressive emphysema, respiratory failure, and in advanced cases the need for long-term oxygen or lung transplantation.
    \item Frequent and severe chest infections.
    \item Cirrhosis of the liver and its complications, including fluid accumulation, bleeding from varices, and hepatic encephalopathy.
    \item An increased risk of hepatocellular carcinoma in those with cirrhosis.
    \item Reduced exercise capacity, disability, and impaired quality of life.
    \item Psychological effects of living with a progressive inherited condition.
\end{itemize}

\section*{Prognosis and Outcomes}
The outlook varies widely and depends strongly on smoking. Non-smokers with AATD may live relatively normal lives with little or no lung disease, whereas smokers typically develop disabling emphysema earlier in life. Liver disease progresses in a minority of people, occasionally to the point of requiring transplantation, but many individuals with AATD never have significant liver problems. Early diagnosis, smoking cessation, vaccination, and pulmonary rehabilitation all improve symptoms and long-term outcomes, and augmentation therapy may slow the decline in lung function in appropriately selected people.

\section*{Prevention and Self-Care}
Because the disease is inherited, prevention focuses on protecting lung and liver health:

\begin{itemize}
    \item Never smoke, and avoid second-hand smoke and dusty or fume-laden environments.
    \item Have recommended vaccinations, including annual influenza and pneumococcal vaccines.
    \item Keep active with regular, gentle exercise and maintain a healthy diet.
    \item Limit alcohol consumption to protect the liver.
    \item Treat chest infections promptly and report any worsening of symptoms.
    \item Attend scheduled lung and liver reviews.
    \item Discuss testing of family members who may also be affected, so they can benefit from early diagnosis.
\end{itemize}

\section*{Sources and Further Reading}
The following reputable sources provide reliable, up-to-date information on alpha-1 antitrypsin deficiency:

\begin{itemize}
    \item NHS. \textit{Alpha-1 antitrypsin deficiency information.} https://www.nhs.uk/
    \item Mayo Clinic. \textit{Alpha-1 antitrypsin deficiency: symptoms, causes, and treatment.} https://www.mayoclinic.org/
    \item MSD Manual. \textit{Alpha-1 antitrypsin deficiency overview.} https://www.msdmanuals.com/
    \item MedlinePlus. \textit{Alpha-1 antitrypsin deficiency health information.} https://medlineplus.gov/
    \item World Health Organization. \textit{Information on genetic disorders and chronic respiratory disease.} https://www.who.int/
    \item NICE. \textit{Clinical guidance on COPD and lung disease.} https://www.nice.org.uk/
\end{itemize}
